% !TEX root = ../main.tex

\begin{digest}

\section{Introduction and Motivation}

Remotely operated vehicles (ROV) have become indispensable platforms for underwater inspection, maintenance, and environmental monitoring in shallow-water environments such as aquaculture facilities, offshore wind farms, and coastal infrastructure. Unlike deep-sea operations, shallow-water zones present distinct challenges: surface wave coupling, near-shore tidal currents, frequent sediment resuspension, multi-path acoustic interference from dual-boundary reflections, and dense artificial structures with flexible obstacles such as ropes and net cages. These environmental factors simultaneously degrade navigation accuracy, visual perception quality, and acoustic communication reliability while amplifying propulsion power demands. Under such conditions, high-precision, low-latency closed-loop motion control becomes critically important.

Compared with autonomous underwater vehicles (AUV), ROV offer decisive advantages in shallow-water missions through tethered real-time communication, continuous power delivery, and the ability to leverage topside computational resources. However, many existing small-ROV platforms rely on open-source autopilot ecosystems---typically Pixhawk flight controllers running ArduSub firmware with pulse-width modulation (PWM) thruster interfaces and QGroundControl ground stations. While such configurations lower the barrier to building prototype vehicles, their general-purpose design limits the depth of customization available for thruster control protocols, device scheduling strategies, and power protection mechanisms. Furthermore, conventional PWM-based thruster interfaces lack error checking and telemetry feedback, restricting closed-loop observability and disturbance rejection capabilities at the actuator level.

This thesis addresses the problem of high-precision motion control for a small, custom-built ROV operating under the practical constraints of limited sensing modalities (no Doppler velocity log, no ultrashort-baseline acoustic positioning, no external motion capture) and modest embedded computational resources (STM32F4 microcontroller). The overarching goal is to establish a complete closed-loop methodology spanning dynamics modeling, control algorithm design, embedded software and hardware development, and quantitative experimental validation---all while maintaining engineering feasibility in a resource-constrained setting.

\section{Dynamics Modeling and Simplification}

A complete six-degree-of-freedom (6-DOF) nonlinear dynamics model is established based on Fossen's vectorial framework, describing the ROV's motion through two coordinate systems: a North-East-Down (NED) earth-fixed frame for absolute position and orientation, and a body-fixed frame aligned with the vehicle's principal axes for velocities, forces, and moments. The full model incorporates rigid-body inertia, added mass, Coriolis and centripetal effects, hydrodynamic damping, and restoring forces arising from gravity and buoyancy.

To make the model suitable for real-time embedded implementation, several engineering simplifications are introduced. Neutral buoyancy is achieved through careful ballast design, eliminating the net vertical restoring force in equilibrium. The center of gravity is positioned directly below and collinear with the center of buoyancy, providing passive roll and pitch stability while eliminating off-diagonal coupling terms in the restoring force vector. The inertia and damping matrices are diagonalized under the assumption of low-speed, near-horizontal operation with three-plane symmetry, retaining dominant diagonal terms while discarding numerically small off-diagonal coupling coefficients. The resulting simplified nominal model preserves the essential dynamic characteristics---including inertial lag, viscous damping, and gravitational restoring moments---while substantially reducing the computational burden for online control calculations.

\section{High-Precision Motion Control Algorithm Design}

The control architecture adopts a hierarchical, decoupled structure that separates attitude regulation from translational velocity control. This design reflects the practical observation that under normal operating conditions, the ROV maintains a near-horizontal attitude, making cross-coupling between roll, pitch, and yaw secondary effects that can be treated as disturbances rather than explicit control targets.

\textbf{SQRT-PID Dual-Loop Attitude Control.} The attitude controller employs a cascaded dual-loop structure. The outer loop maps attitude errors (roll, pitch, yaw) to desired angular velocity commands through a square-root proportional-integral (SQRT-PI) law. The square-root mapping compresses large errors, preventing excessive control effort during aggressive maneuvers, while the integral term with anti-windup protection eliminates steady-state bias. The inner loop tracks the desired angular velocities using a high-frequency PID controller, providing active damping and rapid disturbance rejection. This dual-loop structure effectively decouples the trade-off between transient responsiveness and steady-state accuracy: the outer loop governs the envelope of acceptable angular rates, while the inner loop ensures precise velocity tracking.

\textbf{Velocity Closed-Loop Translation Control.} Surge, sway, and heave directions are regulated through independent velocity controllers. Depth control is achieved by closing a velocity loop on the vertical axis, with the depth error converted to a desired heave velocity that the velocity controller then tracks. This approach simplifies tuning and provides a natural interface for manual piloting and autonomous station-keeping.

\textbf{Model-Based Feedforward and Disturbance Compensation.} To reduce the burden on the feedback controller and improve response speed, known dynamic effects---including rigid-body inertia, added mass, Coriolis forces, hydrodynamic damping, and restoring moments---are pre-computed from the simplified nominal model and applied as feedforward compensation. In parallel, an extended state observer (ESO) estimates unmodeled dynamics, parameter uncertainties, and external environmental disturbances (e.g., currents, wave-induced forces) as an extended state, which is then compensated through an additional feedforward channel. This dual feedforward structure---physics-based plus observer-based---enables the feedback controller to operate with smaller gains while maintaining robust performance.

\textbf{Thrust Allocation for Over-Actuated Configuration.} The ROV is equipped with six thrusters arranged in a non-coplanar configuration that provides full 6-DOF actuation, resulting in an over-actuated system. A thrust allocation matrix is derived from the geometric layout of the thrusters, mapping each thruster's force vector to the resulting generalized forces and moments on the vehicle body. The pseudo-inverse of this allocation matrix yields the minimum-norm (and thus minimum-energy) thruster command vector for any desired control output, distributing the control effort across available actuators while respecting the over-actuated nature of the system.

\section{Embedded Software and Hardware Development}

\textbf{Unified Device Scheduling Framework.} The embedded control software is built around a modular, extensible scheduling framework centered on a unified device base class (\textbf{DeviceBase}) implemented in C++ on the STM32F4 platform. Every hardware peripheral (IMU, depth sensor, thrusters, power monitor) and logical module (attitude controller, command decoder, state estimator) inherits from this base class and implements two standard interfaces: \textbf{Update()} for reading sensor data or refreshing internal state, and \textbf{Handle()} for executing control logic or issuing commands. All device instances are registered into a global queue upon construction. The scheduler iterates through the queue in a defined order: bottom-up \textbf{Update()} calls ensure that low-level sensor data propagates upward before high-level modules consume it, while top-down \textbf{Handle()} calls guarantee that control commands flow from the decision layer down to actuators in a deterministic sequence. Each device supports independent frequency division, allowing IMU sampling at hundreds of hertz while power monitoring runs at a few hertz, all within the same scheduling framework.

\textbf{DSHOT Digital Thruster Protocol.} To overcome the limitations of conventional PWM-based thruster control---including limited resolution (typically 8--12 bits), susceptibility to electromagnetic interference, absence of error detection, and lack of telemetry feedback---the DSHOT digital communication protocol is implemented. DSHOT encodes throttle commands as digital pulses with CRC checksums, enabling reliable transmission at update rates up to several kilohertz with 11--16-bit resolution. The implementation leverages STM32F4's GPIO-DMA hardware acceleration: throttle values are converted to BSRR register mask sequences offline, then transferred to GPIO output registers via DMA at precise clock-triggered intervals, achieving the strict timing requirements of the DSHOT protocol (e.g., 300~kHz for DSHOT300) without CPU polling overhead. After throttle transmission, the GPIO port switches to input mode to sample electronic speed (eRPM) feedback pulses from the electronic speed controllers (ESCs). Sampling follows the Nyquist criterion at three times the DSHOT transmission frequency, and the captured pulse train is decoded and CRC-verified to extract actual propeller speed. This closed-loop speed feedback enables accurate thrust control and provides the foundation for the thrust-to-throttle mapping described in Chapter~5.

\textbf{Power Protection System.} A dedicated power protection board based on the LM5069 hot-swap controller is designed to safeguard the ROV's electrical system. The board provides: (1) soft-start inrush current limiting through a TIMER-pin capacitor ($C_T = 330$~nF, calculated inrush time $\approx 165$~ms), preventing MOSFET damage and battery voltage sag during power-on; (2) under-voltage lockout (UVLO) at approximately 18~V and over-voltage lockout (OVLO) at approximately 32.5~V, set via external resistive dividers, protecting against battery over-discharge and supply over-voltage; and (3) over-current protection through a sense resistor ($R_{\text{SENSE}} = 0.5$~m$\Omega$) with current limiting at approximately 110~A. Experimental validation confirmed the expected soft-start voltage ramp, correct UVLO/OVLO threshold switching, and stable operation under rated loads.

\section{Simulation and Experimental Validation}

\textbf{Hydrodynamic Parameter Identification via CFD.} The added mass matrix and hydrodynamic damping coefficients are essential for accurate dynamics simulation and model-based control. These parameters are obtained through computational fluid dynamics (CFD) simulations in Simcenter STAR-CCM+. Recognizing that added mass (acceleration-dependent inertial forces) and damping (velocity-dependent viscous and pressure drag) are fundamentally coupled in constant-velocity simulations, two separate simulation strategies are employed. For added mass, steady-state simulations with a constant inflow velocity of 0.5~m/s are performed using a steady solver; under fully developed steady flow, acceleration-related terms vanish, allowing the added mass diagonal terms to be extracted from the force equilibrium in each degree of freedom. For damping coefficients, planar motion mechanism (PMM) transient simulations are conducted, in which the ROV undergoes forced single-DOF motions (sinusoidal oscillation or constant-velocity translation) in each of the six degrees of freedom. A $k$-$\omega$ SST turbulence model with incompressible water is used throughout. The simulation domain employs an overset mesh strategy with an inner moving cylindrical region (diameter 1.2~m, length 1.6~m) enclosing the ROV and an outer stationary rectangular domain ($10 \times 8 \times 5$~m). Wall pressure forces, viscous forces, and corresponding moments are recorded for each DOF, and linear and quadratic damping coefficients are extracted via least-squares fitting from the force-velocity data. The resulting diagonal added mass, linear damping, and quadratic damping matrices are integrated into the simulation model and the feedforward compensator.

\textbf{Thruster Characterization and Thrust Lookup Table.} The mapping between ESC throttle commands and actual propeller thrust is affected by supply voltage, propeller characteristics, and motor-ESC dynamics. A dedicated underwater thrust measurement bench is constructed using an aluminum profile frame, a single-axis force sensor, a data acquisition unit communicating via MODBUS-RTU, and the STM32F4 controller issuing DSHOT throttle commands. Thrust is measured across a range of supply voltages and throttle values in both forward and reverse directions, producing a discrete voltage--throttle--thrust dataset. A two-dimensional inverse lookup table is constructed, mapping desired thrust and measured bus voltage to required throttle command. Online interpolation is performed using the ARM CMSIS-DSP library's \textbf{arm\_bilinear\_interp\_f32} function, optimized for the Cortex-M4 floating-point unit, ensuring real-time thrust-to-throttle conversion in the embedded control loop.

\textbf{MATLAB/Simulink Closed-Loop Simulation.} A 6-DOF closed-loop simulation framework is built in MATLAB/Simulink, integrating the simplified dynamics model, CFD-derived hydrodynamic parameters, the SQRT-PID dual-loop controller, model feedforward, ESO disturbance observer, and pseudo-inverse thrust allocation. The simulation framework is structured following the modular approach of the BlueROV2 open-source benchmark simulator, with each subsystem (reference generation, controller, dynamics, thrust allocation, state feedback) encapsulated as an independent block. Position and attitude tracking simulations demonstrate that the closed-loop system achieves stable convergence with no divergence or sustained oscillation. The simulated response exhibits inertial lag and damping characteristics consistent with the CFD-derived hydrodynamic parameters, confirming correct integration of the dynamics model with the control algorithm.

\textbf{Still-Water Field Experiments.} Real-world validation is conducted in a still-water environment. Prior to experiments, the ROV undergoes buoyancy trimming, thruster direction verification, sensor bias calibration, and communication link testing. The embedded controller handles thruster driving, sensor acquisition, and attitude estimation using the onboard extended Kalman filter (EKF) fusion engine, while a topside computer issues desired position, attitude, and motion commands. Depth tracking experiments show that the ROV responds rapidly to target depth commands, with the actual depth converging quickly to the setpoint and maintaining small-amplitude fluctuations around the target. Attitude tracking experiments across roll, pitch, and yaw demonstrate smooth and stable response: the dual-loop SQRT-PID structure prevents excessive control effort under large attitude errors while maintaining high sensitivity to small deviations, and the angular velocity inner loop provides effective damping that eliminates overshoot and oscillation. The experimental results confirm that the integrated control solution---dual-loop attitude control, velocity closed-loop translation, dynamics feedforward, and optimized thrust allocation---achieves fast response, high steady-state accuracy, and stable operation on the physical ROV platform.

\section{Conclusions and Outlook}

This thesis has presented a complete, closed-loop methodology for high-precision ROV motion control in shallow-water environments under resource-constrained conditions. The main contributions are: (1) a systematically simplified 6-DOF dynamics model that preserves essential physical characteristics while meeting embedded real-time constraints; (2) an SQRT-PID dual-loop attitude controller with model-based feedforward and ESO disturbance compensation that balances responsiveness, steady-state accuracy, and robustness; (3) a modular, extensible embedded software framework with unified device scheduling that coordinates multi-peripheral, multi-rate tasks on resource-limited hardware; (4) a DSHOT-based digital thruster interface with hardware-accelerated DMA transmission and closed-loop speed feedback, replacing conventional PWM; (5) a comprehensive LM5069-based power protection solution with soft-start and multi-threshold protection; and (6) systematic validation through CFD parameter identification, thrust bench characterization, Simulink simulation, and still-water field experiments. The results demonstrate that the proposed approach achieves reliable, high-performance motion control suitable for inspection, station-keeping, and close-range intervention tasks in shallow-water environments.

Future work should address: adaptive control strategies for time-varying environmental disturbances (waves, currents); vision-based perception and visual servoing for autonomous intervention; online thruster health monitoring and fault-tolerant thrust allocation; and multi-sensor fusion with SLAM for long-duration autonomous navigation. These extensions would further enhance the operational capability and autonomy of small ROVs in increasingly complex shallow-water missions.

\end{digest}
